\chapter{Report Plot Capture and Zoom}
\label{ch:plots}

\section{Overview}
The design report (Chapter~\ref{ch:report}) is not just tables---it also embeds the analysis
plots: bending-moment, shear-force, and deflection envelopes. This chapter covers two related
pieces of plot work: capturing Matplotlib figures into the report as cropped PNG images
(PR~\#173), and a zoom-window tool for the on-screen plots that works around Matplotlib's lack
of a 3D camera-fit API (PR~\#185).

\section{Capturing Plots into the Report}
To appear in the PDF, a plot must be turned into an image. The capture logic lives in
\texttt{core/bridge\_types/plate\_girder/plot\_generator.py}. The
\texttt{capture\_report\_figures} function builds the envelope plots \emph{fully off-screen} and
returns a dictionary of PNG bytes---one entry per figure---so no GUI or Matplotlib import is
needed on the report side:

\begin{lstlisting}[caption={Off-screen capture of the report figures (\texttt{plot\_generator.py}).},label={lst:capture}]
def capture_report_figures(ds_all, nodes, members,
                           edge_dist=0.0, eng_scale=1.0) -> dict:
    """Build Envelope ULS analysis plots for the report - fully off-screen."""
    data = {}
    data['bm_envelope'] = _render_report_figure(bm_fig, show_max=True,  show_min=True)
    data['sf_envelope'] = _render_report_figure(sf_fig, show_max=True,  show_min=True)
    data['defl_ll']     = _render_report_figure(defl_fig, show_max=True, show_min=False)
    return data
\end{lstlisting}

\subsection{Figure-to-bytes and auto-cropping}
Each figure is rendered by \texttt{\_render\_report\_figure}, which annotates the plot, saves it
to an in-memory PNG buffer, and---importantly---auto-crops the surrounding white border so the
image sits neatly in the report rather than floating in whitespace. The cropping uses PIL to
find the bounding box of the non-white content and adds a few pixels of padding:

\begin{lstlisting}[caption={Rendering to PNG bytes with automatic white-border cropping (\texttt{plot\_generator.py}).},label={lst:crop}]
def _render_report_figure(fig, show_max: bool, show_min: bool):
    """Apply annotations, save to PNG bytes, auto-crop white borders, close."""
    raw_buf = BytesIO()
    fig.savefig(raw_buf, format='png', dpi=150)
    plt.close(fig)

    img = Image.open(raw_buf).convert('RGB')
    bg  = Image.new('RGB', img.size, (255, 255, 255))
    bbox = ImageChops.difference(img, bg).getbbox()   # non-white bounding box
    if bbox:
        left, top, right, bottom = bbox
        img = img.crop((left - 4, top - 4, right + 4, bottom + 4))  # 4px padding
    out = BytesIO(); img.save(out, format='png')
    return out.getvalue()
\end{lstlisting}

The captured bytes are stored directly in the \texttt{figure\_data} dictionary and handed to the
report generator. On the desktop side, \texttt{output\_dock.py} calls
\texttt{capture\_report\_figures} from its report handler and merges the result into
\texttt{figure\_data} before generation. A review point on this PR---that Matplotlib should not
be imported in the UI layer---was addressed by keeping all figure work inside the
\texttt{core} plot generator and passing only bytes across to the frontend.

\section{The Plot Zoom Window}
The 3D CAD view has a true camera and can zoom-to-window through OpenCascade
(Chapter~\ref{ch:toolbar}). The 2D/3D Matplotlib plots have no such API. PR~\#185 solves this by
\emph{physically scaling the canvas} inside a \texttt{QScrollArea} and scrolling to the selected
region---the effect looks identical to a camera zoom-window.

The user drags a rubber-band rectangle over the plot; on release, \texttt{\_execute\_zoom\_window}
computes the scale factor that makes the selected rectangle fill the viewport, resizes the
canvas, and scrolls so the selection is centred:

\begin{lstlisting}[caption={Zoom-to-window by scaling the canvas in a scroll area (\texttt{mpl\_plot\_widget.py}).},label={lst:zoomwin}]
def _execute_zoom_window(self, rect):
    """Zoom so the selected screen rectangle fills the scroll-area viewport."""
    factor    = min(vp_w / rect.width(), vp_h / rect.height())
    new_scale = min(self._zoom_scale * factor, 8.0)   # cap the zoom
    self._zoom_scale = new_scale
    self._apply_zoom()                                # physically resize canvas
    h_bar.setValue(...); v_bar.setValue(...)          # scroll to centre selection
\end{lstlisting}

The rubber-band selection itself is handled by a small set of Matplotlib event callbacks
(\texttt{\_zwin\_on\_press}, \texttt{\_zwin\_on\_motion}, \texttt{\_zwin\_on\_release}), enabled
by \texttt{\_toggle\_zoom\_window}. The toolbar wires the button to the plot through
\texttt{\_plots\_toggle\_zoom\_window} in the controller, keeping the zoom-window mode mutually
exclusive with pan and rotate.

\section{Result}
Analysis plots are captured into the report as tightly-cropped PNG images produced entirely
off-screen, and the on-screen plots gained a zoom-window tool that behaves like the CAD view's
despite Matplotlib having no native equivalent.

\osdagfig{figure8.1.png}{A bending-moment envelope plot captured and auto-cropped for
the design report.}{fig:plotcapture}

\osdagfig{figure8.2.png}{The plot zoom-window: dragging a rectangle scales the canvas so
the selection fills the viewport.}{fig:plotzoom}

\section{Summary}
This chapter added off-screen Matplotlib capture with automatic cropping to feed plots into the
design report, and a canvas-scaling zoom-window that gives the 2D plots a camera-like
zoom-to-region despite Matplotlib's API limitations.
